% !TeX root = ../main.tex
% ============================================================
\section{Discussion}
\label{sec:discussion}
% ============================================================


% ------------------------------------------------------------
\subsection{Geometry-consistent trajectories and dual local adaptation}
% ------------------------------------------------------------

The central methodological result is that recurrent disturbance evidence was more consistently expressed relative to the recent state and variability of each stand than by absolute Sentinel-1 backscatter alone. Absolute VV and VH distributions overlapped substantially, whereas historical disturbances showed highly repeatable within-stand declines. This distinction motivates the three linked elements of the framework: acquisition geometry is kept consistent within each track-specific trajectory, overlapping windows repeatedly renew the local temporal reference, and a window-specific permutation null calibrates each CUSUM amplitude against the variability of that same local context.

Keeping Sentinel-1 tracks separate is important because systematic differences among acquisition geometries would otherwise enter the temporal background against which disturbance is evaluated. Track separation does not eliminate moisture, phenological, or structural variability, but it provides an internally comparable observation sequence within which local adaptation can operate. Windowed CUSUM then addresses the remaining temporal heterogeneity at two levels. Local-reference renewal allows later disturbances to be evaluated against the evolving stand condition, while local statistical calibration means that the same raw CUSUM amplitude need not represent the same level of evidence in a stable and a highly variable trajectory.

The component analysis supports this interpretation. Local-reference windowing produced the largest gain in recurrent-event sensitivity, whereas central-interval admission and permutation calibration progressively recovered precision. The benchmark results showed the same pattern: compared with decrease-filtered ReCuSum, Windowed CUSUM retained similar precision but substantially higher recall, particularly in polygons containing multiple recorded disturbances. The methodological contribution is therefore not windowing alone. CuSum NRT already applies CUSUM to short sliding time series for near-real-time monitoring~\citep{Ygorra2024CuSumNRT}, while ReCuSum handles repeated changes through recursive segmentation~\citep{Ygorra2023ReCuSum}. Here, overlapping local contexts, central event admission, and window-specific statistical calibration are coupled to reconstruct repeated events along a long stand trajectory without making later detections conditional on earlier segmentation.

These comparisons refer to the evaluated stand-level adaptations and should not be interpreted as demonstrating superiority over the complete operational CuSum NRT, CCDC, or published polarimetric BOCD systems.


% ------------------------------------------------------------
\subsection{Recurrent disturbance histories and rotation-relevant interpretation}
% ------------------------------------------------------------

The methodological framework enabled individual detections to be translated into recurrent stand histories. Across the monitored balsa inventory, consecutive detected disturbances showed a dominant approximately three-year timescale, with a median interval of 36 months. This central tendency is consistent with the rapid growth and short commercial development periods reported for balsa in Ecuador~\citep{GonzalezOsorio2010Balsa,CanadasLopez2019BalsaGrowth}, but it should not be interpreted as a directly observed planting-to-harvest age.

% The resulting subset is therefore a robustness check rather than a corrected rotation estimate.
Cross-polarization analysis supports the persistence of this recurrence pattern while also illustrating its uncertainty. Requiring VH support at both bounding events retained a subset with a median interval of 40 months, but VH agreement preferentially retained longer intervals. More generally, balsa rotation length is expected to vary with site productivity, management, and market conditions, and longer commercial cycles have been reported in other production systems~\citep{Jenkin2019BalsaValueChain,Midgley2010Balsa}.

Several features of the detected histories further limit a direct biological interpretation. Planting dates were unavailable, the first and last cycles were censored by the observation period, omitted disturbances can lengthen observed intervals, and additional detections can shorten them. The presence of intervals below 12 months also indicates that not every pair of successive negative transitions corresponds to a complete harvest--regrowth cycle. Individual detections may represent clear-cutting, partial removal, repeated clearing, or other structural changes. The 36-month estimate is therefore best interpreted as a rotation-relevant recurrence timescale of observed stand turnover within the Sentinel-1 record.

The value of the near-decadal archive lies in this ability to observe repeated transitions within the same management units. Longer temporal coverage increases the number of complete adjacent-event intervals that can be reconstructed, extending SAR monitoring beyond isolated disturbance alerts toward multi-cycle plantation histories.


% ------------------------------------------------------------
\subsection{Robustness, limitations, and transferability}
% ------------------------------------------------------------

The evaluation design addressed several sources of optimistic performance estimation but does not eliminate all uncertainty. Spatial blocking kept geographically related polygons, tracks, and events together during parameter selection and evaluation, reducing spatial leakage~\citep{Roberts2017SpatialCV}. However, performance remains conditional on the available reference catalogue and does not constitute independent geographic validation. Unrecorded disturbances can reduce apparent precision, while reference dates partly derived from remote-sensing interpretation may retain some sensor-related dependence.

Statistical calibration is another important boundary. The independent-value permutation provided the highest F1 among the evaluated null formulations at the common threshold, but dependence-preserving alternatives shifted the detector toward higher precision and substantially lower recall. This result should not be interpreted as proof that monthly Sentinel-1 observations are exchangeable. Rather, the retained permutation null represents an empirically effective operating choice under the evaluated precision--recall trade-off. Stronger statistical error-control claims would require independently observed no-disturbance periods, explicit null simulations, or dependence-aware resampling schemes that preserve realistic short-range temporal structure~\citep{Kunsch1989BlockBootstrap}.

Temporal truncation also distinguishes retrospective event reconstruction from immediate alerting. Recall increased substantially as post-disturbance observations accumulated, reaching 0.790 after six months, while 96.8\% of events already detected at that horizon were dated within two months of their corresponding 12-month estimates. Thus, limited post-event support primarily affected whether a disturbance accumulated sufficient evidence for detection, rather than the stability of its estimated month once detected. Monthly compositing further means that estimated event months should not be interpreted as exact felling dates.

Disturbance attribution remains a separate limitation. Negative Sentinel-1 transitions can be associated with harvesting but may also arise from partial canopy removal, flooding, storm damage, or other surface-state changes, and apparent performance depends on reference design and disturbance definition~\citep{RodriguezPaulino2024Review}. The prospective 2026 review was candidate-conditioned and therefore informs candidate precision but cannot estimate missed disturbances or population recall.

The transferable contribution is consequently not a fixed balsa-specific threshold, but a framework in which observation geometry is controlled before detection and both temporal background and statistical evidence are adapted locally. This structure is potentially applicable to other short-rotation plantation systems with repeated management disturbances and limited optical continuity. Future work should determine when provisional candidates have accumulated sufficient post-disturbance evidence to be considered temporally stable and develop dependence-aware statistical calibration that preserves realistic Sentinel-1 temporal structure without the substantial recall loss observed under conservative block-based nulls.